% Generated shared reviewed material.
% Full proofs for actual divisor incidence, signed common-d series and finite laboratory.
\section{Proofs and limits of the signed-divisor development}\label{sd:details}
\subsection{Actual divisor incidence and a finite error certificate}
The identity
$\mathbf1_{(m,n)=1}=\sum_{d\mid m,\ d\mid n}\mu(d)$ is an identity of
ideal divisors. Expanding the finite sums proves~\eqref{sd:incidence}.
When $m=da$ is squarefree, $(d,a)=1$. If $(d,u)=1$, multiplication by
$\chi_d(u)$ is a rowwise unit phase; if $(d,u)>1$ it is the original zero.
For original $x_m=X_U(m)\nu(m)W(Nm/D)$, precisely
$F_d(x;u)=\nu(d)\chi_d(u)T_d^X(u)$ on $(d,u)=1$, and it is zero
otherwise. Thus their row norms agree on the actual restricted support.
For general vectors the coefficient $x_{da}$ remains unfactored.
No identity $X_U(da)=X_U(d)X_U(a)$ is used. Residual $a,b$ may overlap.

Write $\|x_d\|_2^2=\sum_{d\mid m}|x_m|^2$.
The hypotheses of Proposition~\ref{sd:lift} give
\[
 \begin{split}
 |DQ_1(x,y)|&\le\sum_d |Z_d(x,y)|
 \le\sum_d\lambda_d\|x_d\|_2\|y_d\|_2\\
 &\le\left(\sum_d\lambda_d\|x_d\|_2^2\right)^{1/2}
       \left(\sum_d\lambda_d\|y_d\|_2^2\right)^{1/2}
 =\|x\|_B\|y\|_B.
 \end{split}
\]
This is a divisor-incidence lift of positive-row bounds to an indefinite
coprime form. It is not an inference from positivity of a larger Gram
matrix. In particular $B(m)\le\lambda_{\max}\tau_{\rm ideal}(m)$,
with $\tau_{\rm ideal}(m)=2^{\omega(m)}$ on squarefree ideals.
The subpolynomial divisor bound recovers the existing exponent; this
triangle does not supply the desired signed cancellation.

For the original unweighted two-pivot block ($w_u=1$), the accepted
arbitrary-coefficient inputs
at $L=D/Nd\ge1$ allow
\[
 \lambda_d\ll_\epsilon(HD)^\epsilon
 \min\left\{R+\frac{LR}{M}+\frac{L^{2/3}R}{M^{1/3}},\ L+R^2\right\},
 \qquad R=A_1A_2,\quad M=\max(A_1,A_2).
\]
Their constants are classical-input constants, not computed numbers.
Any removal of a $d$-row mask occurs only inside the positive norm.
For $1/2\le L<1$, $Na\in[L,2L]$ permits at most the unit ideal; its
operator norm is bounded directly by the actual row measure.
For an empty subsequence take $\lambda_d=0$.

For the fully computable choice $\lambda_d=E_d$, the matrix entries are
$\sqrt{w_u}\chi_m(u)$ on $d$-divisible columns. The squared Frobenius
norm is precisely the stated raw nonzero-entry count. It is valid for
arbitrary coefficient vectors, not just the saved trajectory. No unknown
sieve constant enters that finite certificate.
For approximations $t_x,t_y$ with errors $e_x=x-t_x,e_y=y-t_y$,
\[
 |Q_1(x,y)-Q_1(t_x,t_y)|
 \le D^{-1}\bigl(\|e_x\|_B\|y\|_B+
                    \|t_x\|_B\|e_y\|_B\bigr).
\]
The AA specialization follows by expanding three terms. To verify it
without floating square roots, put $T=\|t\|_B^2,E=\|e\|_B^2$ and
$L=D|Q_1(x,x)-Q_1(t,t)|-E$.
If $L>0$, the certificate is the equivalent real inequality
$L^2\le4TE$; if $L\le0$ it is immediate. In the saved exact
$\mathbf Z[\omega]$ computations with rational weights, these are
integer/rational checks rather than floating-root tests.
These are finite coefficient-error statements. They do not prove a
varying-family remainder estimate, spectral positivity or a determinant
comparison.

\subsection{A signed common-divisor Euler product}
From here restore ALL original annular squarefree good columns, not an
arbitrary subset from the vector lemma, and use the actual unweighted
raw row block of Section~\ref{sw:core}.
On the stated nonprincipal blocks it is the corresponding component of
the residue-subtracted moment. On a block containing principal rows the
raw identity remains valid, but the original principal correction is
additional and is not discarded.
Expand the exact $d$ identity while keeping its sign. For fixed $a,b,u$
the inner sum is
\begin{equation}\label{sd:inner}
 \sum_{\substack{d\ \mathrm{squarefree}\\(d,Sabu)=1}}
 \mu(d)A_U(da)A_U(db)
       W(Nd\,Na/D)\overline{W(Nd\,Nb/D)}.
\end{equation}
The exact outer coefficient is
$\nu(a)\overline{\nu(b)}\chi_a(u)\overline{\chi_b(u)}$.
The residual $a,b$ are individually squarefree but need not be coprime.
If their quotient character loses a common prime, that prime remains in
$\mathcal T=S\cup\operatorname{supp}(abu)$ and in every original zero.
The common $d$ character cancels only to $\mathbf1_{(d,u)=1}$.

For $\Re s>1$ the series $\mathcal Z$ is absolutely convergent, with
coefficient majorant
$3^{\omega(a)+\omega(b)}9^{\omega(d)}$, uniformly in $U$.
Use four positive Perron variables. Write $\mathcal E_{\mathcal T}(s;v,w)
=\prod_{p\notin\mathcal T}(1-(Np)^{-s}P_pQ_p)$.
The fixed residual prefactor is
$P_a(v)Q_b(w)=\prod_{p\mid a}P_p(v)\prod_{p\mid b}Q_p(w)$;
a common prime of $a,b$ contributes both factors.
The common $d$ prime has nine joint cutoff positions and the sign
$\mu(p)=-1$, giving local factor $1-zPQ$ with
$z=(Np)^{-s}$, $\alpha_i=(Np)^{-v_i}$ and $\beta_j=(Np)^{-w_j}$.
Thus the exact correction in~\eqref{sd:nine} is
\begin{equation}\label{sd:local-correction}
 H_p=\frac{(1-zPQ)\prod_{i=1}^2(1-z\alpha_i)
                           \prod_{j=1}^2(1-z\beta_j)}
                {(1-z)\prod_{i,j=1}^2(1-z\alpha_i\beta_j)}.
\end{equation}
The linear coefficient of numerator and denominator is
$-1-\sum_{i,j}\alpha_i\beta_j$, fixing all nine zeta signs.
At all four auxiliary variables zero the product is
$1/\zeta_K^{\mathcal T}(s)$ and $H_p=1$.
At $v_2=w_2=0$ it is $1/\zeta_K^{\mathcal T}(s+v_1+w_1)$, again
with $H_p=1$. These are generating-function checks, not finite-cutoff
identifications.

For $\Re s\ge1/2+\eta$, $\Re v_i,\Re w_j\ge0$, the degree-five
numerator and denominator of~\eqref{sd:local-correction} have equal
constant and linear coefficients. The sums of absolute coefficient
majorants in degrees two through five are at most $146$ and $26$.
Since $|PQ|\le9$ and $|\alpha_i|,|\beta_j|\le1$,
\begin{equation}\label{sd:normal}
 |H_p-1|\le C_\eta(Np)^{-1-2\eta},\qquad
 C_\eta=\frac{172}{(1-7^{-1/2-\eta})^5}.
\end{equation}
Indeed the five denominator factors have modulus at least $1-|z|$;
good prime ideal norms are at least seven. Ideal counting proves normal
convergence, uniform in deleted primes and vertical phases, and
\[
 |\mathcal H_{\mathcal T}|
 \le\exp\left(C_\eta\sum_p(Np)^{-1-2\eta}\right).
\]
Local correction factors can vanish; no inverse-correction bound is
asserted. Normal convergence controls neither the four numerator nor
five reciprocal zeta factors in the critical region. Their poles, zeros
and crossed residues cannot be suppressed.

\subsection{Finite Perron heights and the exact eight-state numerator}
Set $V=\lfloor U\rfloor+1/2$ and
$\Delta_V=\min_{n\ge1}|\log(V/n)|\ge1/(2V+1)$.
For four finite vertical heights $T_j$ on positive real lines $c_j$,
let $I_T(s)$ be the fourfold integral of
$P_aQ_b\mathcal E_{\mathcal T}(s;v,w)$ against
$\prod_j V^{z_j}\,dz_j/(2\pi i z_j)$.
In the initial half-plane $\sigma=\Re s>1$, finite-height Fubini is
justified by the absolute placement majorant. Put
\[
 \begin{split}
 A_{\rm abs}(\sigma,c)={}&
 \prod_{p\mid a}(1+q^{-c_1}+q^{-c_2})
 \prod_{p\mid b}(1+q^{-c_3}+q^{-c_4})\\
 &{}\cdot\prod_{p\notin\mathcal T}
 \{1+q^{-\sigma}(1+q^{-c_1}+q^{-c_2})
                     (1+q^{-c_3}+q^{-c_4})\},\quad q=Np.
 \end{split}
\]
Then
\begin{equation}\label{sd:perron-error}
 |I_T(s)-\mathcal Z_{U;a,b,\mathcal T}(s)|
 \le V^{\sum c_j}A_{\rm abs}(\sigma,c)
 \left\{\prod_{j=1}^4
    \left(1+\frac{2}{\pi T_j\Delta_V}\right)-1\right\},
\end{equation}
and $A_{\rm abs}\le3^{\omega(a)+\omega(b)}\zeta_K(\sigma)^9$.
For one short integer norm $n$, the indicator is at most $(V/n)^c$,
and the truncated Perron error is at most
$2(V/n)^c/(\pi T\Delta_V)$. Expand the difference of four products and
sum every nonempty error subset to prove~\eqref{sd:perron-error}.
Consequently successive Perron limits recover the exact cutoffs.
The half-integer gap also gives a height-uniform $C_U(V/n)^c$ bound,
allowing dominated convergence in the initial absolutely convergent
series. Four infinite Perron integrals are not declared absolutely
convergent or interchanged silently with Mellin inversion.

There is an entirely finite alternative. Partition
$a=a_0a_1a_2$ and $b=b_0b_1b_2$, internally coprime in each column.
Cross-column overlaps remain allowed. A common $d$ prime occupies one
state $(i,j)\in\{0,1,2\}^2$. Separate its $(0,0)$ part $e$ from eight
pairwise coprime squarefree atomic ideals $h_{ij}$ avoiding $\mathcal T$,
and let $h$ be their product. Define
\[
 h_i^L=\prod_{j=0}^2h_{ij}\ (i=1,2),\qquad
 h_j^R=\prod_{i=0}^2h_{ij}\ (j=1,2).
\]
An atomic $h_{ij}$ with both indices nonzero belongs to one short
factor in each column. These four short factors must not be forced
globally disjoint. Let $\mathcal A_U$ comprise precisely these tuples
satisfying
$N(a_i h_i^L)\le U$ and $N(b_jh_j^R)\le U$ for $i,j=1,2$.
The numerator in~\eqref{sd:finite-numerator} is exactly
\begin{equation}\label{sd:eight}
 \mathcal N_{U;a,b,\mathcal T}(s)=
 \sum_{\mathcal A_U}
 \frac{\mu(a_1a_2)\mu(b_1b_2)
       \mu(h_{11}h_{12}h_{21}h_{22})}{(Nh)^s}
       \prod_{p\mid h}(1-(Np)^{-s})^{-1}.
\end{equation}
At a common prime the sign is
$(-1)^{1+\mathbf1_{i>0}+\mathbf1_{j>0}}$.
It is negative for $(0,0)$, positive for the four single-short states,
and negative for the four double-short states. Summing the remaining
squarefree $e$, coprime to $\mathcal T h$, gives
$1/\zeta_K^{\mathcal T}(s)$ times the local product in~\eqref{sd:eight}.
This proves the exact identity, including its sign.

Multiplying the four cutoff inequalities yields $Nh\le U^4$; double
states are counted twice in that product. Hence the numerator is finite
and holomorphic on $\Re s>0$. At $U=1$ it equals one: this boundary is
not a high-window counterexample or a global pole obstruction.
For $\sigma\ge\sigma_0>0$ and every $\epsilon>0$,
\begin{equation}\label{sd:numerator-bound}
 |\mathcal N_{U;a,b,\mathcal T}(\sigma+i\tau)|
 \ll_{K,\sigma_0,\epsilon}(Na\,Nb)^\epsilon
               U^{4\max(1-\sigma,0)+\epsilon}.
\end{equation}
There are at most $3^{\omega(a)+\omega(b)}$ residual partitions and
$8^{\omega(h)}$ atomic placements at fixed $h$. The finite local product
is $\ll_{\sigma_0,\epsilon}(Nh)^\epsilon$. Apply the ideal-divisor bound
and ideal counting up to $U^4$, using smaller internal exponents; a
logarithm at the transition is absorbed in $U^\epsilon$.
Uniformity as $\sigma_0\to0$ is not claimed. For
$\sigma=1/2+\eta$ with $0<\eta<1/2$, the old-type cutoff cost is
$U^{2-4\eta+\epsilon}$, not a new saving.

\subsection{A uniform common-profile Mellin embedding of the actual mean}
For residual $a,b$ put $M=\max(Na,Nb)$ and
$\varrho=\min(Na,Nb)/M$. At $\varrho\le1/2$ the product of the two
profile values vanishes, including equality because of endpoint
vanishing. Otherwise put
\[
 \Phi_\varrho(t)=W(t)\overline{W(\varrho t)}
\]
when $Na\ge Nb$, and use its conjugate for the reversed orientation.
It is supported in $[1/\varrho,2]\subset[1,2]$.
With $M_k(W)=\max_{j\le k}\|(t\,d/dt)^jW\|_\infty$,
\begin{equation}\label{sd:profile-mellin}
 |\widehat\Phi_\varrho(\sigma+i\tau)|
 \le C_{I,k}M_k(W)^2(1+|\tau|)^{-k}
\end{equation}
uniformly for $\sigma$ in any fixed compact strip $I$ and all
$\varrho\in[1/2,1]$. Leibniz and log-coordinate integration by parts
prove this bound; no ratio-dependent change of $W$ is made.
The transform is entire. For $c>1$, ordinary Mellin inversion after the
finite numerator identity gives the exact scalar~\eqref{sd:inner} as
\begin{equation}\label{sd:mellin}
 \frac1{2\pi i}\int_{\Re s=c}
 \widehat\Phi_\varrho(s)(D/M)^s
 \frac{\mathcal N_{U;a,b,\mathcal T}(s)}
                      {\zeta_K^{\mathcal T}(s)}\,ds.
\end{equation}
This integral is vertically absolutely convergent in its initial
half-plane. The original profile forces $M\le2D$, so the residual $a,b$
and actual row sums are finite.

Keep the full local deletion factor
\[
 \frac1{\zeta_K^{\mathcal T}(s)}=
 \frac1{\zeta_K(s)}\prod_{p\in\mathcal T}(1-(Np)^{-s})^{-1}.
\]
At fixed $\sigma_0>0$ it costs
$\ll_{S,\sigma_0,\epsilon}(N\operatorname{rad}(abu))^\epsilon$,
hence at most $(HD^2)^\epsilon$ on actual support.
This estimates the finite mask dependence, not the central reciprocal.
Multiplying~\eqref{sd:mellin} by every original outer coefficient and
summing the actual rows and residuals gives
\begin{equation}\label{sd:coherent}
 DS_1^{AA}=\frac1{2\pi i}\int_{\Re s=c}
                   \frac{D^s\mathcal A_{D,H,U,W}(s)}{\zeta_K(s)}\,ds.
\end{equation}
Explicitly, with each profile orientation retained,
\[
 \begin{split}
 \mathcal A_{D,H,U,W}(s)
 ={}&\sum_{u\in\mathcal R}
 \sum_{\substack{a,b\ \mathrm{squarefree,good}\\Na,Nb\le2D,\ \varrho>1/2}}
 \nu(a)\overline{\nu(b)}\chi_a(u)\overline{\chi_b(u)}\,M^{-s}\,\widehat\Phi_\varrho(s)\\
 &{}\cdot\mathcal N_{U;a,b,\mathcal T}(s)
       \prod_{p\in\mathcal T}(1-(Np)^{-s})^{-1}.
 \end{split}
\]
This is a \emph{finite coherent sum}, not an independent marginal
moment or a positive sum. All raw zeros remain in its outer coefficients.
For fixed finite $D,H,U$ it is holomorphic on $\Re s>0$ and rapidly
decreasing on fixed vertical strips there, with constants depending on
those finite parameters. No uniform target estimate for it is supplied.
This is an exact signed embedding, not a zero-location theorem.

\subsection{A nonzero short-state correction and the unproved step}
An actual annular control uses good prime ideals of norms
$7,13,19,31,43$ with primary generators
$(1,3),(4,3),(-2,3),(1,6),(7,6)$ in the $1,\omega$ basis.
Take $D=1500,U=25$, $d=p_7$, $a=p_{13}p_{19}$,
$b=p_{13}p_{31}$ and the nonprincipal row $u=p_{43}$ in the pure block
$A_1=37,A_2=1,H=241$, literal class one.
The full raw supports are disjoint from $u$, and
$N(da)=1729,N(db)=2821$ lie in the annulus. A single fixed smooth
profile with plateau $[1.1,1.9]$ equals one at both nodes.
Direct assignments give
$A_U(a)=A_U(b)=-1$, $A_U(da)=1$, $A_U(db)=-1$.
The actual signed $d$ coefficient is $+1$. The $(0,0)$ state alone
would give $-1$; states $(1,0)$ and $(2,0)$ each contribute $+1$.
Dropping all short states therefore misses a nonzero $e=1$ correction.
This illustrates the new eight-state identity, not a new discovery of
coefficient nonmultiplicativity.

The original pair has gcd $p_7p_{13}$. Its contributions across
$d=1,p_7,p_{13},p_7p_{13}$ cancel exactly. Thus a fixed-residual
reciprocal does not prove an uncancelled pole or an obstruction in the
complete coherent mean. The missing step is to estimate that full mean
by $DH(HD)^\epsilon$, retaining fixed-profile/seminorm dependence,
raw masks and conductor/ray/capacity restrictions. Scalar Mertens
cancellation in the long state, a normal correction product, or an
absolute cutoff numerator bound does not prove it. Principal and short
assembly, independent positive determinant/zero-location bridges and
separate Dirichlet, Hecke, Dedekind, elliptic and Artin obligations remain
open. No holomorphy or desired zero location is assumed to close them.

\subsection{Exact arithmetic and source-bound computational scope}
The evaluator streams each actual row and accumulates original
$x_m\chi_m(u)$ into every ideal-divisor bucket $d\mid m$. It then
forms all nine cross products for $A_U,C_U,\mu$ and sums with $\mu(d)$.
Write $n_{\rm col}=|\mathcal C|$,
$I=\sum_m\tau_{\rm ideal}(m)$, $K$ for the distinct divisors and $r$
for the listed rows. This bucket stage costs $O(r(I+K))$ arithmetic operations with a
fixed three channels. Its streaming storage is $O(K)$ plus incidences
and either a streamed row or separately retained phase table.
Character construction, factorization and integer bit lengths are
additional costs, not declared free.
For coefficients bounded by $C$, bucket coordinates are $O(n_{\rm col}C)$,
products $O(n_{\rm col}^2C^2)$ and accumulated values
$O(r n_{\rm col}^2C^2)$ for unit row
weights. For compressed integer multiplicities replace $r$ in this
bit-size bound by $\sum_u w_u$ and account for weight multiplication.

For $D=60000$, the actual phase table has shape $20\times16953$;
the incidence has $73420$ nonzeros over $21750$ divisor buckets, and
$F_d$ has $20\times21750\times3$ complex states.
There are $271688516$ ordered coprime column pairs.
The naive all-pair row loop would visit $5748084180$ pairs and was
\emph{not executed} at this size. This is an operation count, not an
unexecuted timing benchmark. Independent direct pair evaluation covers
the two complete small families; the large family is checked linearly
from literal field arithmetic and divisor-poset inversion.

The exact large results also include
\[
 S_1^{CC}=\frac{661}{7500},\qquad
 S_1^{C\mu}=\frac{13931}{30000},\qquad
 S_1^{\mu\mu}=-\frac{229}{1250}.
\]
For the ordered vectors $(A_U,C_U,\mu)$ the full pairing table is
\[
 \frac1{30000}
 \begin{pmatrix}
 -30714&-11287&19427\\
 -11287&2644&13931\\
 19427&13931&-5496
 \end{pmatrix}.
\]
This is the table of three original coefficient vectors, not a spectral
carrier for GRH. The conjugate $\mu C$ entry and the AA identity retain
their actual linear-first convention. All three saved computations use source v1.
A separate v2 helper fixes its unused $N=U=1$ high-window assertion;
every one of the $17158$ executed columns has $N>U$, so no accepted
row/table or source snapshot was silently replaced.
The fixed-profile flat-node implementation refuses requests above its
stated $D\le60000$ certificate. Neither these three different finite
blocks nor their measured runtimes imply an asymptotic saving.
